%! Semi-log Plots — Unit 2, Lesson 2.15 — Group Activity
\documentclass[10pt]{article}
\usepackage{precalculus-article}
\usepackage{precalculus-boxes}
\newcommand{\TallMath}[1]{$\displaystyle #1\rule[-1.4em]{0pt}{3.2em}$}

\begin{document}

\pageheader{Unit 2, Lesson 2.15}{Group Activity}
\namepartnerperiod

\vspace{0.1in}

\begin{scenariobox}[Drug Concentration Over Time]{sky}
A pharmacologist is studying how a drug concentration $C$ (mg/L) in the bloodstream decays
over time $t$ (hours).  Your group will use semi-log linearization to determine whether an
exponential model is appropriate and, if so, to construct that model.
\end{scenariobox}

\vspace{0.1in}

\begin{tcolorbox}[colback=white, colframe=black!40,
    title=\textbf{Tier R --- Remediate},
    fonttitle=\bfseries, arc=2mm, left=3mm, right=3mm, top=2mm, bottom=2mm]

\textbf{Given data:}
\vspace{4pt}
\begin{center}
\renewcommand{\arraystretch}{1.8}
\begin{tabular}{|c|c|c|c|c|c|}
\hline
\rowcolor{sky}
$t$ (hours) & $0$ & $1$ & $2$ & $3$ & $4$ \\
\hline
$C$ (mg/L) & $100$ & $50$ & $25$ & $12.5$ & $6.25$ \\
\hline
\end{tabular}
\end{center}

\begin{enumerate}[leftmargin=1.8em, itemsep=10pt]

\item Complete the table of $\log_{10}(C)$ values. Round to three decimal places.

\vspace{4pt}
\begin{center}
\renewcommand{\arraystretch}{1.8}
\begin{tabular}{|c|c|c|c|c|c|}
\hline
\rowcolor{sky}
$t$ & $0$ & $1$ & $2$ & $3$ & $4$ \\
\hline
$\log_{10}(C)$ & \blank{1.5cm} & \blank{1.5cm} & \blank{1.5cm} & \blank{1.5cm} & \blank{1.5cm} \\
\hline
\end{tabular}
\end{center}

\item What is the rate of change of $\log_{10}(C)$ per unit of $t$?

\vspace{4pt}
Rate of change $\approx $ \blank{4cm} per hour.

\item Is the rate of change constant?  What does this tell you about the appropriate model
      for this data?

\writelines{2}

\end{enumerate}
\end{tcolorbox}

\vspace{0.1in}

\begin{tcolorbox}[colback=white, colframe=black!40,
    title=\textbf{Tier A --- Approaching Proficiency},
    fonttitle=\bfseries, arc=2mm, left=3mm, right=3mm, top=2mm, bottom=2mm]

\textbf{Continue using the data from Tier R.}

\begin{enumerate}[leftmargin=1.8em, itemsep=10pt]

\item Write the linear model: $\log_{10}(C) = m\cdot t + c$.  Identify $m$ and $c$.

\vspace{4pt}
$m = $ \blank{3cm} \qquad $c = $ \blank{3cm}

\item Find the base $b$ of the exponential model: $b = 10^m \approx $ \blank{2cm}.
      What does the value of $b$ mean in context (i.e., what happens to the drug
      concentration each hour)?

\writeline

\item Find the initial value: $a = 10^c = $ \blank{2cm}.

\item Write the exponential model: $C(t) = $ \blank{5cm}.

      Verify by computing $C(1)$: \blank{4cm}. Does this match the table? \blank{1.5cm}

\item Predict $C(6)$ using the model.  Show your work.

\vspace{4pt}
$C(6) = $ \blank{5cm}

\end{enumerate}
\end{tcolorbox}

\vspace{0.1in}

\begin{tcolorbox}[colback=white, colframe=black!40,
    title=\textbf{Tier E --- Extension},
    fonttitle=\bfseries, arc=2mm, left=3mm, right=3mm, top=2mm, bottom=2mm]

\textbf{New data set --- a different drug trial:}
\vspace{4pt}
\begin{center}
\renewcommand{\arraystretch}{1.8}
\begin{tabular}{|c|c|c|c|c|c|}
\hline
\rowcolor{sky}
$t$ & $0$ & $1$ & $2$ & $3$ & $4$ \\
\hline
$C$ (mg/L) & $80$ & $45$ & $30$ & $22$ & $18$ \\
\hline
\end{tabular}
\end{center}

\begin{enumerate}[leftmargin=1.8em, itemsep=10pt]

\item Compute $\log_{10}(C)$ for each $t$ value. Round to three decimal places.

\vspace{4pt}
\begin{center}
\renewcommand{\arraystretch}{1.8}
\begin{tabular}{|c|c|c|c|c|c|}
\hline
\rowcolor{sky}
$t$ & $0$ & $1$ & $2$ & $3$ & $4$ \\
\hline
$\log_{10}(C)$ & \blank{1.5cm} & \blank{1.5cm} & \blank{1.5cm} & \blank{1.5cm} & \blank{1.5cm} \\
\hline
\end{tabular}
\end{center}

\item Do the $\log_{10}(C)$ values change at a constant rate?  What does this suggest
      about whether an exponential model is appropriate?

\writelines{2}

\item Despite the imperfect fit, compute an approximate slope using the endpoints:
\[
  m \approx \frac{\log_{10}(18) - \log_{10}(80)}{4 - 0} = \frac{\blank{2cm} - \blank{2cm}}{4} \approx \blank{2cm}.
\]

\item Using this slope and the $t = 0$ value as the intercept, write the approximate
      linearized model and recover the corresponding exponential model.

\vspace{4pt}
$\log_{10}(C) \approx $ \blank{4cm}

\vspace{4pt}
$C(t) \approx $ \blank{5cm}

\item Compare this data set to the Tier R data.  Which data set is better described by
      an exponential model?  Justify your answer using evidence from the semi-log
      linearization.

\writelines{3}

\end{enumerate}
\end{tcolorbox}

\end{document}
